\section*{Capítulo \ref{ch:b1:nucleophilic-substitution} --- Sustitución nucleófila}

\begin{solution}{exo:b1:nucleophilic-substitution:1}
\ce{CH3CH2OH + Br-}; \ce{CH3OCH2CH3 + I-}; \ce{CH3CH2CH2CN + Cl-};
\ce{CH3NH3+ + Br-}. \omterm{def:b1:nucleophilic-substitution:substitution}{Sustratos}: los halogenoalcanos; \omterm{def:b1:electronic-effects:nucleophile}{nucleófilos}:
\ce{OH-}, \ce{CH3CH2O-}, \ce{CN-}, \ce{NH3}; \omterm{def:b1:electronic-effects:nucleophile}{grupos salientes}: los iones
haluro.
\end{solution}

\begin{solution}{exo:b1:nucleophilic-substitution:2}
Segundo orden: se triplica; ambas divididas por dos, se divide por 4.
Primer orden respecto de RY solo: no cambia; se divide por dos.
\end{solution}

\begin{solution}{exo:b1:nucleophilic-substitution:3}
SN2 (primario, buen \omterm{def:b1:electronic-effects:nucleophile}{nucleófilo}, aprótico); SN1 (terciario, agua); SN2
(metilo, \omterm{def:b1:electronic-effects:nucleophile}{nucleófilo} fuerte); \omterm{def:b1:nucleophilic-substitution:sn1}{solvólisis} SN1 (secundario, \omterm{def:b1:electronic-effects:nucleophile}{nucleófilo}
neutro débil, \omterm{def:b1:intermolecular-forces-solvents:solvent-classes}{disolvente prótico}), quizá con algo de SN2.
\end{solution}

\begin{solution}{exo:b1:nucleophilic-substitution:4}
El cianuro ocupa el lugar opuesto al bromo y los otros tres grupos se dan
la vuelta. En el producto, el grupo CN tiene la primera prioridad, como
la tenía el Br, y las demás prioridades no cambian: la disposición
invertida tiene el descriptor opuesto, \cip{S}-2-metilbutanonitrilo.
\end{solution}

\begin{solution}{exo:b1:nucleophilic-substitution:5}
1. \omterm{def:b1:electronic-effects:cleavage}{Heterólisis}: flecha del enlace C--Br al Br, que da \ce{(CH3)3C+} y
\ce{Br-} (lenta). 2. Un \omterm{def:b1:lewis-resonance-vsepr:lewis}{par no enlazante} del agua va al carbono
catiónico: \ce{(CH3)3C-OH2+}. 3. Un \omterm{def:b1:lewis-resonance-vsepr:lewis}{par no enlazante} de otra molécula de
agua toma un protón del grupo \ce{OH2+}, y el par O--H vuelve al oxígeno.
Solo la primera etapa, lenta, decide la velocidad; el agua, el
disolvente, está en tal exceso que su concentración no cambia.
\end{solution}

\begin{solution}{exo:b1:nucleophilic-substitution:6}
El carbono que lleva el Br es primario, pero a su lado un grupo
terc-butilo llena el espacio situado detrás del enlace C--Br: el \omterm{def:b1:elementary-steps:energy-profile}{estado
de transición} de la SN2 está muy abarrotado. La SN1 necesitaría un
\omterm{def:b1:electronic-effects:intermediates}{carbocatión} primario, demasiado inestable.
\end{solution}

\begin{solution}{exo:b1:nucleophilic-substitution:7}
El yoduro sustituye al yoduro por SN2, cada vez con inversión: el
\cip{S} se convierte en \cip{R} hasta que ambos están presentes por
igual, una \omterm{def:b1:stereochemistry-in-depth:enantiomers}{mezcla racémica}. Cada sustitución convierte una molécula
\cip{S} en una \cip{R}: la rotación disminuye en el equivalente de dos
moléculas, de modo que la rotación decae el doble de deprisa de lo que se
sustituye el \omterm{def:b1:nucleophilic-substitution:substitution}{sustrato}.
\end{solution}

\begin{solution}{exo:b1:nucleophilic-substitution:8}
SN2: bromometano $>$ 1-bromopropano $>$ 2-bromopropano $>$
2-bromo-2-metilpropano (impedimento estérico). SN1: el orden inverso
(estabilidad del \omterm{def:b1:electronic-effects:intermediates}{carbocatión}), y el bromometano y el 1-bromopropano
apenas reaccionan.
\end{solution}

\begin{solution}{exo:b1:nucleophilic-substitution:9}
$x_{\text{inv}} - x_{\text{ret}} = 0.12$ y $x_{\text{inv}} + x_{\text{ret}} = 1$:
\qty{56}{\percent} invertido y \qty{44}{\percent} con retención. Sobre
todo \omterm{def:b1:nucleophilic-substitution:sn1}{racemización} SN1, con un ligero exceso de inversión: el bromuro que
se va sigue tapando su lado cuando llega el agua.
\end{solution}

\begin{solution}{exo:b1:nucleophilic-substitution:10}
$v = k_1k_2[\mathrm{RY}][\mathrm{Nu}]/(k_{-1}[\mathrm{Y^-}] + k_2[\mathrm{Nu}])$
(\cref{prop:b1:nucleophilic-substitution:sn1-law}). Añadir \ce{Y-}
aumenta el denominador: más \omterm{def:b1:electronic-effects:intermediates}{carbocationes} vuelven al \omterm{def:b1:nucleophilic-substitution:substitution}{sustrato} en lugar de
ser capturados, y la velocidad disminuye. Es despreciable cuando
$k_2[\mathrm{Nu}] \gg k_{-1}[\mathrm{Y^-}]$, como cuando el \omterm{def:b1:electronic-effects:nucleophile}{nucleófilo} es
el disolvente.
\end{solution}

\begin{solution}{exo:b1:nucleophilic-substitution:11}
Concentraciones iguales: $-\mathrm{d}c/\mathrm{d}t = kc^2$,
$1/c = 1/C_0 + kt$,
$t_{1/2} = 1/(kC_0) = \qty{5.0e5}{s}$ (unos seis días). Con un exceso de
hidróxido de cincuenta veces, $[\ce{OH-}] \approx 50C_0$ es constante: $c = C_0
e^{-k't}$ con $k' = 50kC_0 = \qty{1.0e-4}{s^{-1}}$, $t_{1/2} = \ln 2/k' =
\qty{6.9e3}{s}$, unas dos horas.
\end{solution}

\begin{solution}{exo:b1:nucleophilic-substitution:12}
El bromuro tendría que expulsar \ce{OH-}, una \omterm{def:b1:predominant-reaction:strong-acid}{base fuerte} y un mal \omterm{def:b1:electronic-effects:nucleophile}{grupo
saliente}: no hay reacción. En HBr concentrado, el OH se protona a
\ce{-OH2+}, cuyo \omterm{def:b1:electronic-effects:nucleophile}{grupo saliente} es el agua. Alcohol terciario: sale el
agua y da el \omterm{def:b1:electronic-effects:intermediates}{carbocatión}, que captura el \ce{Br-} (SN1). Alcohol
primario: el \ce{Br-} ataca el carbono de \ce{R-OH2+} por detrás mientras
sale el agua (SN2).
\end{solution}

\begin{solution}{pb:b1:nucleophilic-substitution:1}
\textbf{1.} Cada molécula que reacciona da un \ce{H+} y un \ce{Cl-}; la
conductividad es la suma de las contribuciones iónicas, proporcional a la
cantidad de iones.
\textbf{2.} $[\mathrm{RCl}]/[\mathrm{RCl}]_0 = (\kappa_\infty -
\kappa)/\kappa_\infty$.
\textbf{3.} $\ln[\mathrm{RCl}] = \ln[\mathrm{RCl}]_0 - kt$: se representa
$\ln(\kappa_\infty - \kappa)$ frente a $t$.
\textbf{4.} 0.693, 0.603, 0.513, 0.423, 0.333, 0.153, $-0.117$, $-0.387$.
\textbf{5.} Sí, con pendiente de $-0.0180$ por minuto en todo el
intervalo: primer orden.
\textbf{6.} Cualquier orden respecto del agua: su concentración es
constante (degeneración del orden).
\textbf{7.} $k = 0.0180/60 = \qty{3.0e-4}{s^{-1}}$.
\textbf{8.} $t_{1/2} = \ln 2/k = \qty{2.3e3}{s}$, \qty{38.5}{min}.
\textbf{9.} $\ln 10/k = \qty{7.7e3}{s}$, \qty{128}{min}.
\textbf{10.} $\ln(\kappa_\infty - \kappa)$: 0.693 en $t = 0$, 0.109 a
\qty{10}{min}, etc.: pendiente de $-0.0584$ por minuto,
$k = \qty{9.7e-4}{s^{-1}}$.
\textbf{11.} $9.7/3.0 = 3.2$.
\textbf{12.} A \qty{30}{min}: $2.000(1 - e^{-0.0584 \times 30}) = 1.653$;
medido, 1.654.
\textbf{13.} SN1: \omterm{def:b1:nucleophilic-substitution:substitution}{sustrato} terciario, \omterm{def:b1:intermolecular-forces-solvents:solvent-classes}{disolvente prótico}, \omterm{def:b1:electronic-effects:nucleophile}{nucleófilo}
neutro.
\textbf{14.} Como en el \cref{exo:b1:nucleophilic-substitution:5}, con
Cl.
\textbf{15.} La ionización lenta solo implica al \omterm{def:b1:nucleophilic-substitution:substitution}{sustrato}: $v =
k[\mathrm{RCl}]$.
\textbf{16.} El hidróxido solo podría actuar en la etapa rápida de
captura, después de la ionización determinante; la velocidad no depende
de él.
\textbf{17.} El \omterm{def:b1:electronic-effects:intermediates}{carbocatión} es plano y \omterm{def:b1:stereochemistry-in-depth:stereocentre}{aquiral}; el agua se adiciona por
las dos caras por igual: alcohol racémico.
\textbf{18.} Una primera barrera alta (ionización) hasta el \omterm{def:b1:electronic-effects:intermediates}{carbocatión},
en un pozo poco profundo, y después una barrera baja hasta la captura: la
primera etapa es la determinante.
\textbf{19.} Formaría un \omterm{def:b1:electronic-effects:intermediates}{carbocatión} primario (demasiado inestable), y el
agua es un \omterm{def:b1:electronic-effects:nucleophile}{nucleófilo} demasiado débil para una SN2 rápida.
\textbf{20.} $k = A\,e^{-E_a/RT}$.
\textbf{21.} $E_a = R\ln(k_2/k_1)/(1/T_1 - 1/T_2)$.
\textbf{22.} $E_a = 8.314 \times \ln(9.74/3.00)/(1/298.15 - 1/308.15) =
\qty{9.0e4}{J/mol}$.
\textbf{23.} La altura de la barrera de la ionización, la \omterm{def:b1:elementary-steps:rds}{etapa
determinante de la velocidad}.
\textbf{24.} $k = 3.0 \times 10^{-4} \exp\bigl(\frac{90\,000}{8.314}
(\frac{1}{298.15} - \frac{1}{318.15})\bigr) = \qty{2.9e-3}{s^{-1}}$.
\textbf{25.} \textbf{$E_a = \qty{90}{kJ/mol}$}.
\end{solution}
